% !TEX root = ../DoAn.tex
\documentclass[../DoAn.tex]{subfiles}
\begin{document}

Chương này trình bày các giải pháp công nghệ được chọn để hiện thực hóa hệ thống, trả lời
câu hỏi \textit{công nghệ nào, vì sao và đánh đổi gì}. Mỗi mục phân tích vai trò thực tế
trong hệ thống, lý do lựa chọn so với lựa chọn thay thế và hạn chế cần nhận thức.

% =======================================================================
\section{Nền tảng phát triển}
\label{section:3.1}

Hệ thống sử dụng \textbf{Java 21} (LTS, hỗ trợ đến 2031) cho toàn bộ 16 service backend,
tận dụng các cải tiến của nền tảng ở mức phù hợp: Record Classes giảm boilerplate cho một
số DTO; Pattern Matching và switch expression giúp logic điều kiện ngắn gọn hơn. Bản thân
Java 21 cũng mở sẵn khả năng dùng Virtual Threads (Project Loom) để tăng thông lượng I/O,
là hướng tối ưu có thể bật khi triển khai ở quy mô lớn hơn. Toàn
bộ mã nguồn tổ chức theo \textbf{Maven multi-module} với một \texttt{pom.xml} cha quản lý
phiên bản tập trung cho 17 module (16 service + \texttt{common}); module \texttt{common} chứa
\texttt{ApiResponse} wrapper, exception chuẩn hóa và DTO hợp đồng sự kiện Kafka dùng chung,
đảm bảo producer và consumer đồng thuận về cấu trúc thông điệp.

\textbf{Spring Boot 3.5}~\cite{springboot_docs} là framework chính cho 16 service. Bốn lý do lựa chọn:
\textbf{(1) Auto-configuration} --- khai báo dependency, Spring Boot tự cấu hình
\texttt{DataSource}, \texttt{KafkaTemplate}, \texttt{ElasticsearchClient}, loại bỏ phần lớn
cấu hình thủ công; \textbf{(2) Hệ sinh thái đồng bộ} --- Spring Data JPA, Spring Data Redis,
Spring Data Elasticsearch, Spring Kafka, Spring Security và Micrometer Tracing tích hợp liền
mạch; \textbf{(3) Spring Boot Actuator} tự động expose \texttt{/actuator/health} và
\texttt{/actuator/prometheus} làm nền cho observability stack; \textbf{(4) Embedded server}
--- Tomcat cho business service, Netty cho Gateway phản ứng bất đồng bộ, giúp mỗi service
đóng gói thành JAR tự chứa, dễ container hóa. Lựa chọn thay thế Quarkus và Micronaut có
thời gian khởi động nhanh hơn nhờ compile-time DI, nhưng Spring Boot vượt trội về độ hoàn
thiện hệ sinh thái và tài liệu --- yếu tố quan trọng hơn trong thời gian phát triển giới hạn.
Các thư viện hỗ trợ chính: \textbf{Lombok} giảm boilerplate qua \texttt{@Data}/\texttt{@Builder},
\textbf{Jackson} xử lý JSON nhất quán cho REST và Kafka payload, \textbf{Flyway} quản lý
schema migration SQL theo thứ tự, \textbf{SpringDoc OpenAPI} tự động sinh tài liệu API.

\begin{table}[H]
\centering
\caption{Các thư viện và phiên bản chính}
\label{table:dev-tools}
\begin{compacttable}
\begin{tabularx}{\textwidth}{|L{3.5cm}|Y|L{3cm}|}
\hline
\textbf{Mục đích} & \textbf{Thư viện / Công cụ} & \textbf{Phiên bản} \\ \hline
Ngôn ngữ / Runtime   & Java (OpenJDK)            & 21 LTS \\ \hline
Framework chính      & Spring Boot               & 3.5.x \\ \hline
Spring Cloud         & Spring Cloud Dependencies & 2025.0.x \\ \hline
Giảm boilerplate     & Lombok                    & 1.18.x \\ \hline
JSON serialization   & Jackson Databind          & (qua Spring Boot BOM) \\ \hline
Database migration   & Flyway                    & 11.x \\ \hline
API documentation    & SpringDoc OpenAPI         & 2.8.x \\ \hline
\end{tabularx}
\end{compacttable}
\end{table}

% =======================================================================
\section{Hệ sinh thái Spring Cloud}
\label{section:3.2}

\textbf{Spring Cloud Gateway}~\cite{springcloud_docs} là điểm vào duy nhất cho mọi request từ client, xây dựng trên
Spring WebFlux và Reactor Netty với mô hình phi đồng bộ non-blocking, xử lý hàng nghìn kết
nối đồng thời với số thread hạn chế. Mỗi route ánh xạ đường dẫn đến service đích qua Eureka
(\texttt{lb://product-service}) rồi áp dụng chuỗi filter. Hệ thống có hai filter tùy chỉnh:
\textbf{AuthHeaderFilter} thực thi Trusted Subsystem Pattern --- loại bỏ các header
\texttt{X-User-Id/Roles/Email} từ client (ngăn giả mạo), xác minh JWT bằng public key JWKS
từ Keycloak, chèn lại thông tin đã xác thực cho service đích; \textbf{GuestSessionFilter}
tạo và duy trì \texttt{sessionId} qua cookie cho người dùng chưa đăng nhập để bảo toàn giỏ
hàng. Bên cạnh đó, Gateway dùng filter \textbf{RequestRateLimiter} tích hợp sẵn của Spring
Cloud Gateway (thuật toán Redis token bucket) cho \texttt{POST /api/orders}
(10~req/s, burst 20) và \texttt{POST /api/flash-sales/*/purchase} (3~req/s, burst 5). CORS
cấu hình tập trung qua biến môi trường \texttt{CORS\_ALLOWED\_ORIGINS}.

\textbf{Netflix Eureka} là Service Registry trung tâm: service đăng ký địa chỉ khi khởi động
và gửi heartbeat mỗi 30 giây; Gateway truy vấn Eureka để lấy danh sách instance và thực hiện
client-side load balancing Round-Robin. Eureka theo mô hình AP (CAP Theorem), phù hợp với TMĐT
nơi tính sẵn sàng quan trọng hơn nhất quán tuyệt đối; cơ chế self-preservation giữ nguyên
registry thay vì xóa hàng loạt khi heartbeat sụt giảm.

\textbf{Spring Cloud Config Server} quản lý cấu hình tập trung: mỗi service khi khởi động
lấy cấu hình qua HTTP theo tên và profile (\texttt{order-service-docker.yml}...), cho phép
thay đổi tại một điểm mà không cần rebuild image. Trong đồ án dùng native profile (đọc từ
classpath container) thay vì Git backend; hạn chế là cần rebuild Config Server image khi
thay đổi cấu hình.

\textbf{OpenFeign} là HTTP client khai báo: lập trình viên định nghĩa interface với annotation
mô tả endpoint đích, OpenFeign sinh implementation tại runtime và tích hợp Eureka để phân giải
địa chỉ thực tế. Bảy Feign call trong hệ thống: cart-service → product-service (snapshot sản
phẩm khi thêm giỏ); order-service → product-service (snapshot giá), order-service →
voucher-service (đặt chỗ/commit/hoàn voucher), order-service → cart-service (xóa giỏ sau khi
đặt hàng); payment-service → order-service (lấy payment context); review-service →
order-service (xác minh lịch sử mua); và flash-sale-service → order-service (đếm đơn phục vụ
reconciliation). Riêng việc đặt chỗ tồn kho không dùng Feign mà đi qua sự kiện Kafka
\texttt{order-created} trong luồng Saga.

\textbf{Resilience4j} bảo vệ toàn bộ Feign call khỏi cascading failure qua ba cơ chế phối
hợp. \textit{Circuit Breaker} giám sát tỷ lệ lỗi trên sliding window 10 lần: vượt 50\%,
circuit chuyển \texttt{OPEN} (fail-fast, trả fallback ngay); sau 10 giây chuyển
\texttt{HALF\_OPEN} thử lại; nếu thành công về \texttt{CLOSED}. \textit{Retry} thử lại tối
đa 3 lần với delay 500ms cho lỗi có thể phục hồi, chỉ kích hoạt khi circuit \texttt{CLOSED}
hoặc \texttt{HALF\_OPEN}. \textit{Time Limiter} áp timeout 3 giây mỗi call; kết hợp ba cơ
chế, worst-case là $3 \times 3 = 9$ giây trước khi trả fallback, không có request bị chặn
vô thời hạn.

% =======================================================================
\section{Apache Kafka --- Message Broker}
\label{section:3.3}

Apache Kafka~\cite{narkhede2017} là trục giao tiếp bất đồng bộ trung tâm, vận hành ở chế độ \textbf{KRaft}
(Kafka Raft metadata mode) --- không phụ thuộc ZooKeeper, giảm số thành phần hạ tầng và
đơn giản hóa triển khai trong Docker Compose. Hệ thống sử dụng 13 Kafka topic, mỗi topic
phục vụ một luồng nghiệp vụ cụ thể, tổng hợp ở Bảng~\ref{table:kafka-topics}.

\begin{table}[H]
\centering
\caption{Danh sách Kafka topic và luồng dữ liệu tương ứng}
\label{table:kafka-topics}
\begin{compacttable}
\begin{tabularx}{\textwidth}{|L{4.5cm}|L{3.5cm}|Y|}
\hline
\textbf{Topic} & \textbf{Producer} & \textbf{Consumer} \\ \hline
\texttt{user-registered}            & identity-service    & user-service \\ \hline
\texttt{order-created}              & order-service       & inventory-service \\ \hline
\texttt{inventory-updated}          & inventory-service   & order-service \\ \hline
\texttt{inventory-failed}           & inventory-service   & order-service \\ \hline
\texttt{payment-requested}          & order-service       & payment-service \\ \hline
\texttt{payment-success}            & payment-service     & order-service \\ \hline
\texttt{payment-failed}             & payment-service     & order-service \\ \hline
\texttt{order-confirmed}            & order-service       & inventory-service, notification-service \\ \hline
\texttt{order-cancelled}            & order-service       & inventory-service, notification-service \\ \hline
\texttt{flash-sale-order-requested} & flash-sale-service  & order-service \\ \hline
\texttt{product-created}            & product-service     & search-service \\ \hline
\texttt{product-updated}            & product-service     & search-service \\ \hline
\texttt{product-deleted}            & product-service     & search-service \\ \hline
\end{tabularx}
\end{compacttable}
\end{table}

Kafka được chọn thay vì RabbitMQ vì ba ưu thế trong bối cảnh hệ thống này. Thứ nhất, Kafka
lưu trữ thông điệp bền vững trên đĩa với retention có thể cấu hình --- OutboxPoller có thể
đọc lại sự kiện từ quá khứ nếu cần tái xử lý hoặc debug sau sự cố. Thứ hai, cơ chế
partition và consumer group đảm bảo thứ tự sự kiện trong cùng partition trong khi vẫn cho
phép xử lý song song --- yếu tố quan trọng để tuần tự hóa sự kiện của cùng một đơn hàng.
Thứ ba, throughput cao nhờ sequential disk I/O và zero-copy transfer, phù hợp với lưu lượng
lớn trong sự kiện flash sale.

Hạn chế của Kafka so với RabbitMQ: độ trễ end-to-end cao hơn do cơ chế batch và polling,
kém phù hợp cho các tương tác cần phản hồi tức thì --- lý do các call đồng bộ trong hệ
thống vẫn dùng OpenFeign. Kafka cũng phức tạp hơn về cấu hình và vận hành, tiêu tốn nhiều
tài nguyên hơn trên môi trường single-host.

% =======================================================================
\section{Hạ tầng dữ liệu}
\label{section:3.4}

\subsection{PostgreSQL 17}

PostgreSQL~\cite{postgresql_docs} là RDBMS chính, phục vụ 10 trong 13 business service theo nguyên tắc Database per
Service. Trên một instance PostgreSQL duy nhất, mỗi service sở hữu một logical database riêng
biệt (\texttt{order\_db}, \texttt{payment\_db}, \texttt{product\_db}...), đảm bảo cách ly dữ
liệu và cho phép mỗi service quản lý schema độc lập thông qua Flyway migration.

PostgreSQL được chọn thay vì MySQL vì bốn lý do cụ thể. \textbf{UUID native}: PostgreSQL hỗ
trợ kiểu \texttt{uuid} làm khóa chính cho hầu hết bảng trong hệ thống --- phù hợp với môi
trường phân tán nơi ID phải duy nhất toàn cầu mà không cần phối hợp trung tâm.
\textbf{JSONB}: một số bảng (ví dụ snapshot thông tin sản phẩm trong đơn hàng) lưu trữ dữ
liệu bán cấu trúc dưới dạng JSONB, cho phép truy vấn JSON với index hiệu quả. \textbf{FOR
UPDATE SKIP LOCKED}: mệnh đề này --- cốt lõi của cơ chế OutboxPoller cho phép nhiều instance
xử lý song song mà không tranh chấp --- được PostgreSQL hỗ trợ đầy đủ. \textbf{MVCC}:
Multi-Version Concurrency Control xử lý tốt tình huống đọc-ghi đồng thời cao trong các
service như order và payment mà không cần khóa đọc.

Hạn chế: 10 database chạy trên một instance PostgreSQL single-host, không có replication hay
failover --- điểm cần cải thiện khi chuyển sang Kubernetes với managed database.

\subsection{Redis 8}

Redis đảm nhiệm bốn vai trò với đặc điểm truy cập khác nhau, tận dụng mô hình single-threaded
và thao tác in-memory với độ trễ dưới mili-giây:

\textbf{(1) Flash sale atomic operations}: flash-sale-service sử dụng Redis làm bộ đếm tồn
kho (\texttt{flash-sale:\{id\}:stock}) và tập hợp người mua (\texttt{flash-sale:\{id\}:buyers}),
khai thác Lua Script để thực thi kiểm tra và giảm slot một cách nguyên tử trong một bước duy
nhất. Đây là vai trò đặc trưng nhất và quan trọng nhất của Redis trong hệ thống.

\textbf{(2) Distributed Lock}: flash-sale-service sử dụng lệnh \texttt{SET NX EX} để tranh
giành quyền thực thi các scheduler (CampaignScheduler, ReconciliationScheduler) khi chạy
nhiều instance, đảm bảo chỉ một instance xử lý mỗi thời điểm và tránh deadlock vĩnh viễn
qua TTL tự động.

\textbf{(3) Cart storage}: cart-service lưu giỏ hàng với key \texttt{cart:\{userId\}} hoặc
\texttt{cart:\{sessionId\}}, tận dụng TTL để tự động dọn dẹp giỏ hàng bỏ hoang mà không
cần cron job.

\textbf{(4) Rate Limiting}: API Gateway lưu trạng thái token bucket trên Redis, định danh
theo \texttt{X-User-Id} cho các endpoint đã xác thực (đặt hàng, flash sale) và theo địa chỉ
IP cho endpoint công khai, cho phép rate limiting hoạt động nhất quán khi Gateway chạy nhiều
instance.

So sánh lựa chọn thay thế: Memcached thiếu Lua scripting và cấu trúc dữ liệu phong phú
(Sorted Set, Set, Hash) nên không đáp ứng được vai trò flash sale; Hazelcast in-process cache
loại bỏ network overhead nhưng phức tạp hơn khi đồng bộ trạng thái giữa instance và tiêu
tốn heap JVM của chính service.

\subsection{Elasticsearch 8}

Elasticsearch~\cite{elasticsearch_docs} phục vụ read model trong CQRS-lite của search-service. Index \texttt{products}
lưu bản sao denormalized của thông tin sản phẩm, được tối ưu cho full-text search đa ngôn
ngữ, lọc kết hợp nhiều tiêu chí (category, brand, price range) và xếp hạng mức độ liên quan.
Đồng bộ từ PostgreSQL qua Kafka: mỗi sự kiện \texttt{product-created/updated/deleted} từ
product-service được search-service tiêu thụ và ánh xạ vào Elasticsearch document tương ứng,
chấp nhận độ trễ đồng bộ ngắn để đổi lấy hiệu năng tìm kiếm.

So với PostgreSQL full-text search (\texttt{tsvector}): Elasticsearch vượt trội về tốc độ
trên tập dữ liệu lớn, hỗ trợ tìm kiếm mờ (fuzzy search), gợi ý tự động (autocomplete) và
phân tích tần suất từ khóa --- những tính năng quan trọng cho trải nghiệm tìm kiếm sản phẩm
trong TMĐT. So với Apache Solr: Elasticsearch có REST API hiện đại hơn và tích hợp tốt với
Spring Data Elasticsearch.

Hạn chế: Elasticsearch tiêu tốn nhiều RAM (JVM heap mặc định 512MB--1GB), cần cấu hình
\texttt{ES\_JAVA\_OPTS} cẩn thận trong môi trường single-host để không ảnh hưởng đến các
container khác.

% =======================================================================
\section{Keycloak --- Bảo mật và quản lý định danh}
\label{section:3.5}

Keycloak 26~\cite{keycloak_docs} được chọn làm Identity Provider (IdP) tập trung, chịu trách nhiệm quản lý người
dùng, xác thực và cấp phát token theo chuẩn OAuth 2.0~\cite{rfc6749} và OpenID Connect. Với Keycloak, hệ
thống không cần tự triển khai logic quản lý mật khẩu, phiên làm việc hay làm mới token ---
tất cả được Keycloak xử lý theo chuẩn đã được kiểm chứng rộng rãi.

Keycloak đảm nhiệm ba chức năng cụ thể. \textbf{Cấp phát token}: sau khi người dùng xác
thực thành công, Keycloak cấp ba loại token --- Access Token (JWT RS256, thời hạn 5 phút),
Refresh Token (thời hạn 30 phút) và ID Token. Access Token được ký bằng private key của
Keycloak; public key được công bố tại JWKS endpoint
(\texttt{/realms/ecommerce/protocol/openid-connect/certs}) để API Gateway xác minh cục bộ
mà không cần gọi lại Keycloak cho mỗi request. \textbf{Làm mới token}: khi Access Token
hết hạn, client dùng Refresh Token để nhận Access Token mới mà không cần đăng nhập lại.
\textbf{Quản trị người dùng}: quản trị viên tạo, sửa, khóa tài khoản và quản lý vai trò
(\texttt{ROLE\_USER}, \texttt{ROLE\_ADMIN}) qua Keycloak Admin Console.

So sánh lựa chọn: tự xây dựng hệ thống xác thực (Spring Security + database người dùng) thiếu
các tính năng bảo mật đã kiểm chứng như brute-force detection, MFA, social login và token
refresh an toàn; Auth0 và Okta là dịch vụ thương mại, phụ thuộc bên thứ ba và giới hạn số
người dùng ở free tier; Keycloak là mã nguồn mở, self-hosted, không giới hạn người dùng và
phù hợp với môi trường học thuật.

Hạn chế: Keycloak là ứng dụng Java nặng, tiêu tốn đáng kể tài nguyên RAM và CPU trên
single-host --- được quản lý bằng \texttt{mem\_limit} trong Docker Compose và chấp nhận như
đánh đổi để có giải pháp bảo mật toàn diện.

% =======================================================================
\section{Next.js --- Tầng giao diện}
\label{section:3.6}

Next.js 16~\cite{nextjs_docs} (App Router, React 19, TypeScript, Tailwind CSS) phục vụ hai nhóm màn hình:
Storefront cho khách hàng và Admin Panel cho quản
trị viên. Access Token lưu trong \texttt{httpOnly cookie} và chỉ gắn vào header
\texttt{Authorization} khi server-side code gọi API Gateway (pattern BFF --- Backend for
Frontend), ngăn JavaScript phía client truy cập token và giảm nguy cơ XSS. Admin Panel
(\texttt{/admin/*}) được bảo vệ bởi middleware kiểm tra cookie và vai trò \texttt{ROLE\_ADMIN};
các mutation dùng Server Actions để token không xuất hiện trong JavaScript trên trình duyệt.

App Router của Next.js 16 cho phép kết hợp linh hoạt Server Components (render HTML phía
server, giảm JavaScript bundle) và Client Components (tương tác động phía client). Storefront
tận dụng Server Components cho các trang catalog, chi tiết sản phẩm và kết quả tìm kiếm ---
cải thiện SEO và thời gian tải lần đầu. Các tính năng realtime như cập nhật giỏ hàng và
trạng thái flash sale dùng Client Components với React state. Next.js được chọn thay vì React
SPA thuần vì hỗ trợ SSR/SSG, routing tích hợp và không cần reverse proxy riêng cho frontend.

% =======================================================================
\section{Docker --- Đóng gói và triển khai}
\label{section:3.7}

Docker~\cite{docker_docs} đóng gói toàn bộ hệ thống --- 16 service backend, 1 frontend Next.js cùng các thành
phần hạ tầng (PostgreSQL, Redis, Kafka, Elasticsearch, Keycloak, Prometheus, Grafana, Zipkin,
Mailpit) --- thành container độc lập, đảm bảo môi trường đồng nhất giữa phát triển và
production. Docker Compose điều phối các container với healthcheck, \texttt{depends\_on} và
network nội bộ.

Hai file cấu hình tách biệt môi trường: \texttt{docker-compose.yml} cho phát triển local
(build từ source, expose đầy đủ cổng) và \texttt{docker-compose.prod.yml} cho production-like
(kéo image từ GitHub Container Registry, đặt \texttt{mem\_limit} và \texttt{mem\_reservation}
phù hợp tài nguyên EC2). Caddy đóng vai trò reverse proxy trước API Gateway, tự động cấp và
gia hạn chứng chỉ TLS qua Let's Encrypt. Docker Compose được chọn thay vì Kubernetes vì phù
hợp với mô hình single-host của đồ án; hạn chế rõ ràng là không có auto-scaling, self-healing
hay rolling update --- các tính năng cần Kubernetes khi hệ thống ra production thật.

% =======================================================================
\section{Các công nghệ và công cụ hỗ trợ}
\label{section:3.8}

\textbf{VNPAY}~\cite{vnpay_docs} là cổng thanh toán tích hợp, hỗ trợ thẻ ATM nội địa, thẻ tín dụng quốc tế
và QR code, cung cấp môi trường sandbox miễn phí cho phát triển. Payment-service ký tham số
bằng HMAC-SHA512 để tạo URL thanh toán; sau khi người dùng hoàn tất, VNPAY gọi lại qua cả
Return URL (cập nhật giao diện) và IPN endpoint (ghi nhận kết quả phía server), đảm bảo kết
quả không bị mất khi người dùng đóng trình duyệt sớm. Phạm vi đồ án giới hạn ở sandbox.

\textbf{GitHub Actions} tự động build Docker image và publish lên GitHub Container Registry
(GHCR) mỗi khi push lên nhánh \texttt{main}; \textbf{AWS EC2} single-host kéo image mới và
khởi động lại container qua \texttt{docker-compose.prod.yml}, kiểm chứng khả năng triển khai
liên tục trong phạm vi đồ án.

\textbf{Prometheus}~\cite{prometheus_docs} thu thập metrics định kỳ 15 giây từ \texttt{/actuator/prometheus} của
từng service (JVM, request rate/latency/error, Kafka consumer lag, Circuit Breaker state);
\textbf{Grafana} trực quan hóa qua ba dashboard provision tự động: \texttt{jvm-overview},
\texttt{spring-boot-overview} và \texttt{ecommerce-saga-overview} (trạng thái đơn hàng
theo thời gian). \textbf{Zipkin}~\cite{zipkin_docs} theo dõi hành trình request xuyên service: Micrometer
Tracing sinh \texttt{TraceId}/\texttt{SpanId} và propagate qua HTTP header \texttt{traceparent}
(W3C) và Kafka header, cho phép tái tạo toàn bộ chuỗi xử lý khi điều tra sự cố.

\textbf{Mailpit} là SMTP server giả lập cho môi trường phát triển, bắt giữ toàn bộ email từ
notification-service và hiển thị qua giao diện web tại \texttt{http://localhost:8025}; thay
thế bằng SMTP thật trên production qua biến môi trường mà không cần thay đổi mã nguồn.

\textbf{JUnit 5 + Mockito} kiểm thử đơn vị, mock Feign client/Kafka template/repository để
kiểm thử logic nghiệp vụ cô lập. \textbf{Testcontainers} khởi tạo container PostgreSQL thật
trong test thay vì H2 in-memory --- cần thiết vì H2 không hỗ trợ \texttt{uuid}, \texttt{jsonb}
và \texttt{FOR UPDATE SKIP LOCKED}; dùng chủ yếu cho OutboxPoller và Idempotent Consumer.
\textbf{k6}~\cite{k6_docs} kiểm thử hiệu năng với bốn kịch bản: smoke test (9/9 pass), catalog stress test
(p95~=~61\,ms, 200~VU), checkout load test (100\% success, 50~VU) và flash sale spike test
(500~VU đồng thời, 0 over-sell) --- kết quả dùng làm bằng chứng định lượng trong
Chương~\ref{chapter:SolutionAndContribution}.

\end{document}
